\documentclass[11pt,a4paper]{article}
\usepackage[margin=2.4cm]{geometry}
\usepackage{amsmath,amssymb,bm,mathtools}
\usepackage{booktabs}
\usepackage{array}
\usepackage{framed}
\usepackage{xcolor}
\usepackage{enumitem}
\usepackage[colorlinks=true,linkcolor=blue,citecolor=blue,urlcolor=blue]{hyperref}

% --- vectors and tensors (Al-Attar et al. 2024 notation, with I for the background inertia tensor)
\newcommand{\bO}{\bm{\Omega}}
\newcommand{\bw}{\bm{\omega}}
\newcommand{\bx}{\mathbf{x}}
\newcommand{\bu}{\mathbf{u}}
\newcommand{\ba}{\mathbf{a}}
\newcommand{\bb}{\mathbf{b}}
\newcommand{\bt}{\mathbf{t}}
\newcommand{\bj}{\mathbf{j}}
\newcommand{\bc}{\mathbf{c}}
\newcommand{\bk}{\mathbf{k}}
\newcommand{\bh}{\mathbf{h}}
\newcommand{\bq}{\mathbf{q}}
\newcommand{\bl}{\bm{\ell}}
\newcommand{\bchi}{\bm{\chi}}
\newcommand{\bI}{\mathbf{I}}
\newcommand{\It}{\tilde{\mathbf{I}}}
\newcommand{\bN}{\mathbf{N}}
\newcommand{\bGam}{\bm{\Gamma}}
\newcommand{\bn}{\hat{\mathbf{n}}}
\newcommand{\be}{\hat{\mathbf{e}}}
\newcommand{\E}{\mathsf{E}}
\newcommand{\Id}{\mathsf{1}}
\newcommand{\R}{\mathsf{R}}
\newcommand{\Sm}{\mathsf{S}}
\newcommand{\K}{\mathsf{K}}
\newcommand{\W}{\mathsf{W}}
\newcommand{\dd}{\,\mathrm{d}}
\newcommand{\pM}{\partial M}
\newcommand{\se}{\sigma_{\mathrm e}}
\newcommand{\scw}{\sigma_{\mathrm{CW}}}
\newcommand{\Aw}{A_{\mathrm w}}

\definecolor{rembar}{RGB}{70,110,170}
\renewenvironment{leftbar}{%
  \def\FrameCommand{{\color{rembar}\vrule width 2.5pt}\hspace{8pt}}%
  \MakeFramed{\advance\hsize-\width\FrameRestore}}%
  {\endMakeFramed}
\newenvironment{remark}[1][Remark]{\begin{leftbar}\small\noindent\textbf{#1.}\ }{\end{leftbar}}

\title{Dynamic rotational feedbacks in the elastic sea-level problem:\\
governing equations and iterative forward solution}
\author{Working notes (prepared with Claude) --- to be checked}
\date{29 September 2026}

\begin{document}
\maketitle

\begin{abstract}
\noindent
This note is the single reference for the dynamic extension of the rotational sea-level theory of
Al-Attar et al.\ (2024; hereafter AA24). It supersedes \texttt{rotation\_alattar\_notation.pdf} and
\texttt{rotation\_forward\_problem.pdf} (hereafter Note~1 and the companion note), whose logic is
audited in Appendix~\ref{app:audit}. The earth model is elastic, the ocean is hydrostatic, there is
no relative angular momentum apart from the leading-order effect of the fluid core, and no external
torque acts. With these assumptions, AA24's static Euler relation between the rotation vector and
the inertia perturbation is replaced by the linearised Liouville equation, and nothing else
changes. The only memory in the problem is then the equatorial rotation vector, a single 2-vector.
The main result is a practical one (\S\ref{sec:solve}): the forward problem is solved by marching in time,
running the existing static fixed-point iteration at each step with one line changed. An exact
superposition method, which needs no change to \texttt{pyslfp}, provides an independent check.
Iterating over the whole time history at once is shown not to work. Extensions (relative angular
momentum, torques, anelasticity) are indicated where they would enter, and the generalised
problem and its reciprocity theorem are given in \S\ref{sec:recip}. Section~\ref{sec:data} shows
that for data corrected for the pole tide in the standard way, the load-driven rotation cancels from
the observables altogether. The dynamic theory then matters mainly for uncorrected data, for polar
motion itself, and for understanding.
\end{abstract}

\tableofcontents

%=======================================================================
\section{Scope, assumptions and notation}
\label{sec:scope}
%=======================================================================

\paragraph{The problem.} Given a direct load $\zeta(\bx,t)$ on $t\in[0,T]$ (for ice,
$\zeta=\rho_i(1-C)\Delta I$), find the sea-level change $\xi$, displacement $\bu$, gravitational
potential perturbation $\phi$ and rotation-vector perturbation $\bw$ as functions of time.

\paragraph{Assumptions.}
\begin{enumerate}[label=(A\arabic*),leftmargin=3em]
  \item \textbf{Elastic, quasi-static deformation.} The momentum equation is split by projecting it
  onto rigid rotations and onto deformations. The rigid-rotation projection, taken in the Tisserand
  frame, is the Liouville equation, in which all inertial terms are kept (\S\ref{sec:liouville}). In
  particular, the Euler force $-\rho\dot\bw\times\bx$ does no work on Tisserand-frame deformations,
  since $\int_M\rho\,(\dot\bw\times\bx)\cdot\bu\dd V=\dot\bw\cdot\int_M\rho\,\bx\times\bu\dd V=0$. It
  therefore enters only through its torque, which is the retained term $\bI\dot\bw$. The deformational projection is the
  quasi-static elastic problem: $(\bu,\phi)$ respond instantaneously to the current load and
  centrifugal potential. The inertial and Coriolis forces dropped there are smaller than the elastic
  restoring force by $(\nu/\omega_0)^2$ and $2(\nu/\Omega)(\Omega/\omega_0)^2\approx3\times10^{-3}\,\nu/\Omega$,
  where $\omega_0$ is the frequency of the gravest normal mode ${}_0S_2$.
  \item \textbf{Hydrostatic ocean, fixed shorelines.} The sea surface is at every instant an
  equipotential of the gravitational plus centrifugal potential, and the ocean function $C$ does
  not change. The ocean-free problem is the special case $C=0$.
  \item \textbf{No relative angular momentum}, $\bh=\mathbf{0}$, except for the fluid core.
  Displacements are measured in a Tisserand frame, with no rigid rotation of the solid Earth.
  Currents, the atmosphere and the transport of meltwater carry no angular momentum relative to that
  frame. The fluid outer core is different: its relative angular momentum $\bh_{\rm core}$ is driven
  by the wobble itself. It is small, but its gyroscopic term $\Omega\E\bh_{\rm core}$ is a fraction
  $A_f/I_1\approx0.11$ of the retained $I_1\dot\bw$. With the core, the frame is the Tisserand frame of
  the mantle, and the Euler-force argument of (A1) integrates over the mantle. Its leading-order effect is included through
  the wobble inertia $\Aw$ (\S\ref{sec:liouville}). Setting $\Aw=I_1$ recovers the strict
  assumption $\bh=\mathbf{0}$.
  \item \textbf{No external torque}, $\bGam=\mathbf{0}$.
  \item \textbf{Reference state as in AA24.} The elastic--gravitational operator comes from a
  spherically symmetric model, and $\bI$ is the observed inertia tensor with $\bO=\Omega\be_3$
  along the axis of greatest moment. Derivations are for general $I_1\le I_2<I_3$. The
  implementation, like \texttt{pyslfp}, uses $I_1=I_2$.
\end{enumerate}
(A3) and (A4) are relaxed in \S\ref{sec:ext}, which shows that any other $\bh$, and $\bGam$, enter
only as additional excitation.

\paragraph{Ordering.} Let $\nu$ be a typical forcing frequency. The theory treats
$\nu/\se=O(1)$ exactly, where $\se=\Omega(I_3-I_1)/\Aw$ is the Euler frequency, and drops terms
of relative size $\epsilon:=\nu/\Omega$ and $\delta:=(I_3-I_1)/I_1\approx3.3\times10^{-3}$.
Section~\ref{sec:order} lists the neglected effects by their actual size. The largest, the fluid
core (about 11\% of the retained inertial term), is included through $\Aw$. The recommended default
configuration is $\Aw=A_m$, an equilibrium ocean and damping with $Q\approx100$.

\paragraph{Notation.} The notation follows AA24, with three changes. The background inertia tensor
is $\bI$ (AA24's $\mathbf{C}$), so the ocean function keeps the symbol $C$ without a clash.
Sea-level change is $\xi$ (AA24's $\Delta SL$), as in Heathcote et al.\ (2026). The gyroscopic matrix is
$\bN$.

\begin{center}
\small
\begin{tabular}{@{}lll@{}}
\toprule
Symbol & Meaning & AA24\\
\midrule
$\bI$, $I_1\le I_2<I_3$ & background inertia tensor and principal moments & $\mathbf{C}$, $C_i$\\
$\bc$ & first-order perturbation of the inertia tensor & $\mathbf{c}$\\
$\bj=\bc\bO$ & inertia perturbation contracted with $\bO$ & $\mathbf{j}$\\
$\It=\mathrm{diag}(I_3-I_1,\,I_3-I_2,\,-I_3)$ & static rotation matrix & $\tilde{\mathbf{C}}$\\
$\bO=\Omega\be_3$, $\bw$ & reference rotation vector and its perturbation & same\\
$\Psi$, $\psi$ & centrifugal potential and its perturbation & same\\
$\Phi$, $\phi$ & gravitational potential and its perturbation (Eulerian) & same\\
$\bu$ & displacement & same\\
$\xi$ & sea-level change & $\Delta SL$\\
$C$ & ocean function & $C$\\
$\sigma$, $\zeta$ & total and direct surface loads & same\\
$\bk$ & AA24 rotational forcing ($\bk=-\bh$ statically) & $\mathbf{k}$\\
$\bh$, $\bGam$ & relative angular momentum, external torque & ---\\
$\bN$ & gyroscopic matrix, eq.~\eqref{eq:N} & ---\\
$\Aw$; $A_f$, $A_m=I_1-A_f$ & wobble inertia; fluid-core and mantle equatorial moments & ---\\
$\E=\begin{pmatrix}0&-1\\1&0\end{pmatrix}$ & $+90^\circ$ rotation of equatorial 2-vectors & ---\\
\bottomrule
\end{tabular}
\end{center}
Subscripts $\perp$ and $\parallel$ denote components orthogonal and parallel to $\be_3$. Where
no confusion arises, $\bw$ and $\bj$ also denote their equatorial 2-vectors $(\omega_1,\omega_2)$ and
$(j_1,j_2)$. Sign conventions are those of AA24. $\Phi$ increases outwards, so
$\nabla\Phi=g\bn$ at the surface of a spherical model, and the geoid perturbation is $-\phi/g$.

%=======================================================================
\section{What the observations see}
\label{sec:data}
%=======================================================================

Before building any rotational machinery, it is worth asking which rotational signals survive in the
data.

\paragraph{Decomposition.} The problem is linear. Any observable can therefore be written as
\begin{equation}
  d=d^\zeta[\zeta]+\mathsf{X}_d\,\bw_\perp,
  \label{eq:dsplit}
\end{equation}
where $d^\zeta$ is the response to the load with the rotation held at zero (the ocean still
responds), and $\mathsf{X}_d$ is the full response to a unit prescribed rotation (\S\ref{sec:schemeB}).
This holds for the static and dynamic theories alike: they differ only in $\bw_\perp(t)$.

\paragraph{Pole-tide corrections.} Standard processing of GRACE (IERS Conventions; Petit \& Luzum 2010),
of satellite altimetry (Desai et al.\ 2015) and of GNSS removes a modelled pole tide. It is computed
from the \emph{observed} polar motion $\bw_{\rm obs}$ relative to a conventional mean pole
$\bar{\mathbf m}(t)$:
$d_{\rm corr}=d-\mathsf{X}^{\rm corr}_d(\bw_{\rm obs}-\bar{\mathbf m})$. Whatever its cause,
$\bw_{\rm obs}$ is the true rotation. If the correction model matches the physics,
$\mathsf{X}^{\rm corr}_d=\mathsf{X}_d$, then (noting that IERS $y_p$ is positive towards $90^\circ$W, so
$\bw_{\rm obs}/\Omega=(x_p,-y_p)$, and that the observed pole differs from $\bw$ by an $O(\epsilon)$
term)
\begin{equation}
  \boxed{\;d_{\rm corr}=d^\zeta[\zeta]+\mathsf{X}_d\,\bar{\mathbf m}\;}
  \label{eq:dcorr}
\end{equation}
The load-driven rotation cancels exactly, whether it is static or dynamic.

\paragraph{Consequences.}
\begin{enumerate}
  \item For pole-tide-corrected products, the observation operator for $\zeta$ is the
  \emph{non-rotating} one, plus a known offset from the conventional mean pole. For trend
  inversions, only $\dot{\bar{\mathbf m}}$ enters. Neither the static nor the dynamic rotational
  theory is then needed for these data.
  \item The rotational physics left in corrected data is the model mismatch
  $(\mathsf{X}_d-\mathsf{X}^{\rm corr}_d)(\bw_{\rm obs}-\bar{\mathbf m})$. Sources include the
  anelastic $k_2$ used by the IERS, the ocean pole-tide model and the mean-pole convention. This
  mismatch acts mainly on the atmospherically and oceanically excited Chandler and annual wobbles,
  not on the load-driven part.
  \item The dynamic theory is needed for data not corrected for the pole tide (usually tide gauges),
  for polar motion itself used as data, and for understanding. For uncorrected data, the true
  rotation also contains excitation that (A3) excludes and the model cannot reproduce.
  \item The altimetry convention used in \texttt{pyslfp} (Lickley et al.\ 2018; Heathcote et
  al.\ 2026) removes only the direct term $\psi/g$ from modelled sea-surface height. It keeps the
  load-driven solid-Earth and ocean pole tides, and so matches neither uncorrected nor corrected data.
\end{enumerate}
All of this depends on the conventions of each product: which mean pole is used, and which
pole-tide models. These should be settled for the specific data sets before Scheme~1
(\S\ref{sec:schemeA}) is used in inversions.

%=======================================================================
\section{The static problem (AA24 \S\S2--3, restated)}
\label{sec:static}
%=======================================================================

\paragraph{Elastic loading.} For a surface load $\sigma$, $(\bu,\phi)$ satisfy
\begin{equation}
  A(\bu,\phi\,|\,\bu',\phi')+\int_{\pM}(\bu'\cdot\nabla\Phi+\phi')\,\sigma\dd S=0
  \qquad\text{for all }(\bu',\phi'),
  \label{eq:weak0}
\end{equation}
where the bilinear form $A$ is symmetric under exchange of the pairs $(\bu,\phi)$ and
$(\bu',\phi')$, and vanishes on linearised rigid-body motions (AA24 eqs.~1--4).

\paragraph{Sea level and load.} With rotation,
\begin{equation}
  \xi=-\frac1g\,(\bu\cdot\nabla\Phi+\phi+\psi)+\frac{\Phi_g}{g},\qquad
  \sigma=\zeta+\rho_w C\,\xi,\qquad \int_{\pM}\sigma\dd S=0,
  \label{eq:sl}
\end{equation}
where the constant $\Phi_g$ is fixed by the mass constraint (AA24 eqs.~10--11, 53).

\paragraph{Centrifugal potential.} From $\Psi=-\tfrac12\|\bO\times\bx\|^2$,
\begin{equation}
  \psi=-(\bO\times\bx)\cdot(\bw\times\bx),\qquad
  \nabla\psi=\bO\times(\bw\times\bx)+\bw\times(\bO\times\bx).
  \label{eq:psi}
\end{equation}
With $\bO=\Omega\be_3$, $\psi=\Omega\,[\,x_3(\omega_1x_1+\omega_2x_2)-\omega_3(x_1^2+x_2^2)\,]$.
The equatorial components of $\bw$ generate a degree-2, order-1 potential. $\omega_3$ generates
degree-0 and degree-2, order-0 terms.

\paragraph{Inertia identity.} The inertia perturbation produced by a load $\sigma$ and the
resulting displacement $\bu$ satisfies, for every constant vector $\bw'$ with associated
$\psi'$ (AA24 eqs.~46--49),
\begin{equation}
  \bw'\cdot\bj+\int_M\rho\,\bu\cdot\nabla\psi'\dd V+\int_{\pM}\sigma\,\psi'\dd S=0 .
  \label{eq:inertia}
\end{equation}
This is a statement about the mass distribution at one instant. It therefore holds unchanged, at
each $t$, in the dynamic problem.

\paragraph{Static rotation.} The steady Euler equation $\bO\times(\bI\bO)=\mathbf{0}$, perturbed
to first order, together with conservation of $\|\mathbf{L}\|$, gives $\It\bw=\bj$ (AA24
eqs.~36--42). The static problem in weak form is
\begin{equation}
  A(\bu,\phi\,|\,\bu',\phi')+\int_{\pM}(\bu'\cdot\nabla\Phi+\phi')\,\sigma\dd S
  +\int_M\rho\,\bu'\cdot\nabla\psi\dd V+(\It\bw-\bj)\cdot\bw'=0,
  \label{eq:weakstatic}
\end{equation}
for all $(\bu',\phi',\bw')$. Here $-\rho\nabla\psi$ is the centrifugal body force (AA24 eq.~50).
AA24's generalised forcings $(\bt,\zeta_\phi,\bk)$ are deferred to \S\ref{sec:gen}. \S\ref{sec:liouville}
shows that both halves of $\It\bw=\bj$ are limits of the Liouville equation.

%=======================================================================
\section{Rotational dynamics: the linearised Liouville equation}
\label{sec:liouville}
%=======================================================================

\subsection{The full first-order equation}

In the Earth-fixed frame, which rotates with $\bO+\bw$, the angular momentum
$\mathbf{L}=(\bI+\bc)(\bO+\bw)+\bh$ obeys
$\dot{\mathbf{L}}+(\bO+\bw)\times\mathbf{L}=\bGam$. Linearising, with $\bI\bO=I_3\bO$, gives
\begin{equation}
  \frac{\dd}{\dd t}\big(\bI\bw+\bj+\bh\big)+\bO\times\big(\bI\bw+\bj+\bh\big)+\bw\times(\bI\bO)=\bGam .
  \label{eq:liouville}
\end{equation}
In components with general $I_1,I_2$:
\begin{align}
  I_1\dot\omega_1+(I_3-I_2)\Omega\,\omega_2&=\Omega(j_2+h_2)-(\dot j_1+\dot h_1)+\Gamma_1,\notag\\
  I_2\dot\omega_2-(I_3-I_1)\Omega\,\omega_1&=-\Omega(j_1+h_1)-(\dot j_2+\dot h_2)+\Gamma_2,
  \label{eq:components}\\
  I_3\dot\omega_3&=-(\dot j_3+\dot h_3)+\Gamma_3.\notag
\end{align}
With $j_\perp=\Omega(c_{13},c_{23})$ these are the classical Liouville equations (Munk \& MacDonald
1960). The steady, unperturbed case is $\bO\times(\bI\bO)=\mathbf{0}$.

For $I_1=I_2$, the equatorial pair takes the compact 2-vector form
\begin{equation}
  I_1\dot\bw+\dot\bj+\dot\bh=\Omega\,\E\big[(I_3-I_1)\bw-\bj-\bh\big]+\bGam ,
  \label{eq:eqE}
\end{equation}
which can be rearranged exactly as
\begin{equation}
  \boxed{\;\dot\bw=\se\,\E\,(\bw-\bchi),\qquad
  \bchi=\frac{1}{I_3-I_1}\Big[\bj+\bh+\frac{1}{\Omega}\E\big(\bGam-\dot\bj-\dot\bh\big)\Big]\;}
  \label{eq:chi}
\end{equation}
with $\se=\Omega(I_3-I_1)/I_1$ for the moment.
Here $\bchi$ is the (geodetic) excitation. Eqs.~\eqref{eq:eqE} and \eqref{eq:chi} were checked
symbolically against eq.~\eqref{eq:liouville}.

\paragraph{Fluid core and the wobble inertia.} A fluid outer core is coupled to the mantle through
pressure on the elliptical core--mantle boundary. To leading order in $\nu/\Omega$, its relative
rotation is $\E\dot\bw/\Omega$, which is $O(\epsilon)$ relative to $\bw$: the core almost follows the
wobble. Its relative angular momentum in the mantle frame is
$\bh_{\rm core}\approx(A_f/\Omega)\E\dot\bw$, where $A_f$ is its equatorial moment. The term
$-\Omega\E\bh_{\rm core}=A_f\dot\bw$ in eq.~\eqref{eq:eqE} is then of the same kind as $I_1\dot\bw$
and a fraction $A_f/I_1$ of it. It cannot be ordered away, and $\dot\bh_{\rm core}$ is smaller by
a further factor $\nu/\Omega$. The net effect is to replace $I_1$ by the mantle moment
$A_m=I_1-A_f$ in the $\dot\bw$ term only (Smith \& Dahlen 1981). Write $\Aw$ for this
\emph{wobble inertia}: $\Aw=I_1$ under the strict form of (A3), and $\Aw=A_m$ with the core. PREM
gives $A_f/I_1=0.113$, so $A_m/I_1=0.887$. From here on,
\begin{equation}
  \se:=\frac{\Omega(I_3-I_1)}{\Aw},
  \label{eq:se}
\end{equation}
with period 304.5\,d for $\Aw=I_1$ and 270.1\,d for $\Aw=A_m$. The static limit does not involve
$\Aw$ and is unchanged. The Love numbers relevant to the wobble also differ slightly when the core is
decoupled (Smith \& Dahlen 1981). That effect is neglected here.

\subsection{Ordering, and the leading-order equation}
\label{sec:order}

Scale time by the Euler frequency, $\tau=\se t$. With $\bh=\bGam=\mathbf{0}$ (the core now being
carried by $\Aw$), eq.~\eqref{eq:chi} becomes
\begin{equation}
  \frac{\dd\bw}{\dd\tau}=\E(\bw-\bchi),\qquad
  \bchi=\frac{1}{I_3-I_1}\Big[\bj-\delta_w\,\E\frac{\dd\bj}{\dd\tau}\Big],\qquad
  \delta_w:=\frac{\se}{\Omega}=\frac{I_3-I_1}{\Aw},
  \label{eq:scaled}
\end{equation}
with $\delta_w=3.3\times10^{-3}$ for $\Aw=I_1$ and $3.7\times10^{-3}$ for $\Aw=A_m$. In the natural
limit $\nu/\se=O(1)$, the rotation responds on the forcing timescale. For $\Aw=I_1$, $\nu/\se$ is
10.2 at a period of one month, 0.83 at one year and 0.08 at ten years. The only small parameter left
is $\delta_w$. The $\dd\bj/\dd\tau$ term is a uniform correction to the excitation of relative size
$\delta_w\,\nu/\se=\nu/\Omega=\epsilon$: about 3\% at monthly periods and 0.3\% at annual periods. It
is dropped, as in the standard geodetic excitation functions (Barnes et al.\ 1983; Gross 2015).
Keeping it would be harmless but not systematically better, because the relative angular momentum
of the transported meltwater, excluded by (A3), is of the same order.

The \textbf{leading-order equations} are
\begin{equation}
  \boxed{\;\dot\bw_\perp=\se\,\E\,(\bw_\perp-\bq),\qquad \bq:=\frac{\bj_\perp}{I_3-I_1};
  \qquad\qquad I_3\,\omega_3=-j_3\;}
  \label{eq:leading}
\end{equation}
with $\se$ from eq.~\eqref{eq:se}.
The axial relation follows from the third of eqs.~\eqref{eq:components} integrated from a
quiescent start. It is algebraic and carries no memory, which is why AA24's use of $\|\mathbf{L}\|$
conservation is exact at this order. (On sub-decadal timescales the core is also largely
decoupled axially, which would replace $I_3$ by the mantle's polar moment. This is not pursued
here, since $\omega_3$ is neglected in the implementation.)

\paragraph{Neglected effects, by size.}
\begin{itemize}
  \item \emph{Anelasticity.} It lengthens the Chandler period by about 8\,d (2\%) and gives
  $Q\approx100$ (Smith \& Dahlen 1981). It can be represented by the damping of \S\ref{sec:timestep}.
  \item \emph{Non-equilibrium ocean pole tide.} Its size is set by ocean dynamics, not by $\nu/\Omega$.
  The observed pole tide is close to the self-consistent equilibrium response (Desai 2002).
  \item \emph{$\dot\bj$ and the meltwater $\bh$:} $O(\epsilon)$, as above.
  \item \emph{Ellipticity of the elastic response}, i.e.\ spherically symmetric Love numbers:
  $O(\delta)$, since the flattening $1/298\approx\delta$. This is comparable to the $O(\epsilon)$
  terms near the annual band ($\epsilon=\delta$ at periods of about 300\,d).
  \item \emph{Solid-Earth inertial and Coriolis forces in the deformation}:
  $(\nu/\omega_0)^2$ and $3\times10^{-3}\epsilon$ (A1). These are negligible.
\end{itemize}
If $\Aw=I_1$ is used, the fluid core ($\approx11\%$ of the retained term) heads this list.

\paragraph{What controls the response.} The dimensionless forcing frequency that matters for the
response is $\nu/\scw$, not $\nu/\se$, where $\scw$ is the free (Chandler) frequency of
\S\ref{sec:spectral}. At one year $\nu/\scw$ is 1.30 ($\Aw=I_1$) or 1.16 ($\Aw=A_m$). The annual
band therefore lies just beyond the resonance, which is why the annual factors in \S\ref{sec:behaviour}
are large and negative.

\paragraph{Static limit.} Setting $\dot\bw=\mathbf{0}$ gives $\bw_\perp=\bq$, i.e.\
$(I_3-I_1)\bw_\perp=\bj_\perp$, which is AA24's eq.~38. The static and dynamic theories share the
same $\bq$. The only difference is that the static theory \emph{assigns} $\bw_\perp=\bq$, whereas the
dynamic theory makes $\bw_\perp$ \emph{relax towards} $\bq$ by circling it (\S\ref{sec:behaviour}).

\paragraph{Matrix form.} For general $I_1,I_2$, eqs.~\eqref{eq:components} at leading order are
\begin{equation}
  \It\bw+\bN\dot\bw=\bj,\qquad
  \bN=\frac{1}{\Omega}\begin{pmatrix}0&-A_{w,2}&0\\ A_{w,1}&0&0\\0&0&0\end{pmatrix},
  \label{eq:N}
\end{equation}
where $A_{w,i}=I_i$, or the mantle moments with the core. Without the core,
$\bN\mathbf{v}=\Omega^{-2}\bO\times(\bI\mathbf{v})$. $\bN$ has no axial part, so the axial row is
exactly $-I_3\omega_3=j_3$. $\bN$ is antisymmetric if and only if $A_{w,1}=A_{w,2}$, in which case
its equatorial block is $(\Aw/\Omega)\E$. The antisymmetry
matters only for reciprocity (\S\ref{sec:recip}). The forward solver works equally well for
$I_1\neq I_2$.

%=======================================================================
\section{The dynamic forward problem}
\label{sec:problem}
%=======================================================================

Collecting eqs.~\eqref{eq:sl}, \eqref{eq:inertia}, \eqref{eq:weakstatic} and \eqref{eq:N}: for
each $t\in[0,T]$ and all test functions $(\bu',\phi',\bw')$,
\begin{equation}
  \boxed{\;A(\bu,\phi\,|\,\bu',\phi')+\int_{\pM}(\bu'\cdot\nabla\Phi+\phi')\,\sigma\dd S
  +\int_M\rho\,\bu'\cdot\nabla\psi\dd V+(\It\bw+\bN\dot\bw-\bj)\cdot\bw'=0\;}
  \label{eq:weakdyn}
\end{equation}
with $\bw(0)=\bw_0$, where $\sigma=\zeta+\rho_wC\xi$, $\xi$ is given by eq.~\eqref{eq:sl}, $\int_{\pM}\sigma\dd S=0$, and
$\bj$ is the inertia perturbation of the current configuration.

\paragraph{Structure.} Everything in eq.~\eqref{eq:weakdyn} except $\bN\dot\bw$ is instantaneous.
If $\bw(t)$ were known, the fields at each $t$ would follow from a static fingerprint problem with
the rotation vector \emph{prescribed}, i.e.\ with no rotation update. The problem therefore splits
into
\begin{enumerate}[label=(\roman*)]
  \item an instantaneous map $(\zeta(t),\bw(t))\mapsto(\xi,\bu,\phi,\bq)(t)$, which is the static
  solver with $\bw$ held fixed; and
  \item the 2-vector ODE $\dot\bw_\perp=\se\E(\bw_\perp-\bq)$, plus the algebraic $\omega_3$.
\end{enumerate}
The whole memory of the system is $\bw_\perp(t)$. Every method in \S\ref{sec:solve} is a way of
coupling (i) and (ii).

\paragraph{Initial state.} Take $t=0$ as the reference epoch, with $\zeta(0)=0$ and $\bw_0=\mathbf{0}$.
If a load $\zeta_0$ has been in place long enough for equilibrium, instead set $\bw_0$ to the
static solution for $\zeta_0$, which also gives $\bq_0=\bw_0$. Starting with $\zeta(0)\neq0$ and
$\bw_0=\mathbf{0}$ is inconsistent: it excites a spurious free wobble whose amplitude equals the
static response to $\zeta_0$. A free wobble genuinely present at $t=0$ is not determined by the
loads, and $\bw_0$ must then be prescribed.

%=======================================================================
\section{Spherically symmetric implementation}
\label{sec:spectral}
%=======================================================================

\subsection{Love numbers and the rotational terms}

With $\sigma=\sum\sigma_{lm}Y_{lm}$ (orthonormal real harmonics, no Condon--Shortley phase), AA24
eq.~60 gives $u_{lm}=h_l\sigma_{lm}$ and $\phi_{lm}=k_l\sigma_{lm}$. Here AA24's $k_l$ includes the
potential of the load itself. The centrifugal perturbation $\psi$ acts on the Earth like a tidal potential. The deformation it
causes adds $h^T_2\psi_{21}$ to the $(2,1)$ coefficient of the displacement and $k^T_2\psi_{21}$ to that
of $\phi$, where $h^T_2$ and $k^T_2$ are tidal Love numbers. Here $\phi$ is the potential of the Earth's
redistributed mass only and does not contain $\psi$ itself. This differs from AA24's loading $k_l$,
which includes the potential of the load. Writing
$\bm\phi_{21}=(\phi^c_{21},\phi^s_{21})$ at radius $b$ (the mean sea-floor radius in
\texttt{pyslfp}):
\begin{equation}
  \bm\psi_{21}=r\,\bw_\perp,\quad r=\sqrt{\tfrac{4\pi}{15}}\,\Omega b^2;\qquad\qquad
  \bq=\beta\,\bm\phi_{21},\quad
  \beta=\sqrt{\tfrac{5}{12\pi}}\,\frac{\Omega b^3}{G(I_3-I_1)} .
  \label{eq:rbeta}
\end{equation}
The first relation follows from $\psi=\Omega x_3(\omega_1x_1+\omega_2x_2)$ and
$Y^c_{21}=\sqrt{15/4\pi}\,x_1x_3/r^2$. The second is MacCullagh's formula: the exterior degree-2
potential of the perturbed mass has $\phi^c_{21}=\sqrt{12\pi/5}\,Gc_{13}/b^3$, and similarly for
$\phi^s_{21}$ with $c_{23}$, while $\bj_\perp=\Omega(c_{13},c_{23})$. Here $\bm\phi_{21}$ is the potential of
\emph{all} perturbed mass (load, internal deformation, and the tidal response to $\psi$) but not
$\psi$ itself. It therefore contains both terms of eq.~\eqref{eq:inertia} automatically. In
\texttt{pyslfp}, $r$ and $\beta$ are \texttt{rotation\_factor} and \texttt{inertia\_factor}, and the static update is
\texttt{angular\_velocity\_change = i * potential\_change\_lm.coeffs[:, 2, 1]}, i.e.\ $\bw_\perp\leftarrow\bq$.

\subsection{The rotational feedback matrix}
\label{sec:K}

Because map (i) is linear,
\begin{equation}
  \bq(t)=\bq^\zeta(t)+\K\,\bw_\perp(t).
  \label{eq:K}
\end{equation}
Here $\bq^\zeta$ is the excitation from the direct load, with the ocean responding and $\bw\equiv\mathbf{0}$.
$\K$ is a constant $2\times2$ matrix giving the excitation raised by a unit rotation, through the
solid Earth's tidal yielding and the equilibrium ocean pole tide. With no ocean,
$\K=\kappa\Id$ with
\begin{equation}
  \kappa=\beta r k^T_2=\frac{k_2^T}{k_s},\qquad k_s=\frac{3G(I_3-I_1)}{\Omega^2b^5}.
  \label{eq:kappa}
\end{equation}

\paragraph{General form of $\K$.} A unit rotation raises $\bm\psi_{21}=r\be_\alpha$. That changes
the degree-2, order-1 inertia by two routes. The solid Earth yields tidally, adding $k^T_2\bm\psi_{21}$
to $\bm\phi_{21}$; this is the $\kappa\Id$ above. The ocean follows the rotational equipotential, and the
degree-2, order-1 part of the resulting water load adds $k_2\bm\sigma_{21}$. The forcing sea level is
$-\gamma_2\psi/g$ plus a constant, where $\gamma_2=1+k^T_2-h^T_2\approx0.69$ is the usual diminishing
factor (dimensionless Love numbers). Hence
\begin{equation}
  \K=\kappa\Id+\K_{\rm oc},\qquad
  \K_{\rm oc}=\frac{1}{k_s}\,\frac{3\rho_w}{5\bar\rho}\,(1+k'_2)\,\gamma_2\;\mathsf{Q},\qquad
  \mathsf{Q}_{\alpha\beta}=\int C\,Y^\alpha_{21}\,\hat\xi^\beta\dd\Omega ,
  \label{eq:Kform}
\end{equation}
where $k'_2$ is the degree-2 load Love number, $\bar\rho$ the mean density, and the integral runs over
the unit sphere. Here $\hat\xi^\beta$ is the sea level produced by a unit forcing $Y^\beta_{21}$, with
the ocean's self-attraction and loading and with mass conserved. Without self-attraction and loading,
$\mathsf{Q}_{\alpha\beta}=\int CY^\alpha Y^\beta\dd\Omega-(\int CY^\alpha\dd\Omega)(\int CY^\beta\dd\Omega)/A_{\rm oc}$,
which measures how much the two harmonics overlap over the ocean. The second term is the
mass-conservation correction. $\K_{\rm oc}$ is positive definite and anisotropic because the ocean
geometry is. For present-day oceans this overlap matrix is
$\begin{psmallmatrix}0.785&-0.021\\-0.021&0.621\end{psmallmatrix}$, and the prefactor is $0.055$.
Together they reproduce the off-diagonal of $\K_{\rm oc}$ in eq.~\eqref{eq:Knum} below almost
exactly. On the diagonal, self-attraction and loading amplify it by 19\%.

\paragraph{$\K$ is symmetric.} Let $(\bu^a,\phi^a,\sigma^a)$ and $(\bu^b,\phi^b,\sigma^b)$ be the
responses to prescribed rotations $\bw^a,\bw^b$ with no direct load. Using the static weak form
without its rotation row, once with each solution as the test function, and using the symmetry of $A$:
\begin{equation*}
  \int_{\pM}(\bu^b\!\cdot\!\nabla\Phi+\phi^b)\sigma^a\dd S+\int_M\rho\,\bu^b\!\cdot\!\nabla\psi^a\dd V
  =\int_{\pM}(\bu^a\!\cdot\!\nabla\Phi+\phi^a)\sigma^b\dd S+\int_M\rho\,\bu^a\!\cdot\!\nabla\psi^b\dd V .
\end{equation*}
Eq.~\eqref{eq:inertia} then gives
\begin{equation*}
  \bw^b\cdot\bj^a-\bw^a\cdot\bj^b
  =\int_{\pM}(\bu^a\!\cdot\!\nabla\Phi+\phi^a+\psi^a)\sigma^b\dd S-\int_{\pM}(\bu^b\!\cdot\!\nabla\Phi+\phi^b+\psi^b)\sigma^a\dd S .
\end{equation*}
Substitute eq.~\eqref{eq:sl} with $\sigma=\rho_wC\xi$. The constants $\Phi_g$ drop out by mass conservation,
and what remains is $-g\rho_w\int C\xi^a\xi^b+g\rho_w\int C\xi^b\xi^a=0$. Hence
$\bw^b\cdot\bj^a=\bw^a\cdot\bj^b$, and with $\bj^a=(I_3-I_1)\K\bw^a$, $\K=\K^T$. For
$I_1\neq I_2$ the same argument shows that the matrix taking $\bw_\perp$ to $\bj_\perp$ is symmetric.

\paragraph{Numbers.} The \texttt{pyslfp} constants give $k_s=0.951$ and, with $k^T_2=0.2986$,
$\kappa=0.314$. \texttt{pyslfp} takes $\Omega=2\pi/86400$\,s$^{-1}$ (solar day). The sidereal
$7.2921\times10^{-5}$\,s$^{-1}$ should be used instead. It lowers $k_s$ by 0.55\%, raises $\kappa$
and $\K$ by the same factor, and shortens the Euler period by 0.27\%, but it changes the Chandler
periods below by less than 0.3\,d. The numbers here use the \texttt{pyslfp} constants. Two static \texttt{pyslfp} solves at $l_{\max}=64$ with
present-day oceans (the method is in \S\ref{sec:schemeB}) give
\begin{equation}
  \K=\begin{pmatrix}0.3655&-0.0012\\-0.0012&0.3550\end{pmatrix},
  \label{eq:Knum}
\end{equation}
which is symmetric to $2\times10^{-15}$. The ocean pole tide therefore adds about $0.04$--$0.05$ to
the solid-Earth $\kappa$, anisotropically.

\subsection{Free wobble}

Substituting eq.~\eqref{eq:K} into eq.~\eqref{eq:leading},
\begin{equation}
  \dot\bw_\perp=\se\,\E\big[(\Id-\K)\bw_\perp-\bq^\zeta\big].
  \label{eq:reduced}
\end{equation}
Because $\K$ is symmetric, $\E(\Id-\K)$ has zero trace. Its eigenvalues therefore satisfy
$\lambda^2=-\det(\Id-\K)$, and because $\Id-\K$ is positive definite they are purely imaginary,
$\pm i\scw$, with
\begin{equation}
  \scw=\se\sqrt{\det(\Id-\K)}.
  \label{eq:scw}
\end{equation}
\begin{center}
\small
\begin{tabular}{@{}lcc@{}}
\toprule
$T_{\rm CW}$ & no ocean & equilibrium ocean\\
\midrule
$\Aw=I_1$ (no core) & 444.0\,d & 476.0\,d\\
$\Aw=A_m$ (fluid core) & 393.7\,d & 422.2\,d\\
\bottomrule
\end{tabular}
\end{center}
The equilibrium ocean lengthens the period by about 30\,d (Dahlen 1976), and the core shortens it by
11\%. The free motion is undamped and, because $\K$ is slightly anisotropic, slightly elliptical.
Smith \& Dahlen (1981) obtained 426.7 sidereal days (425.5\,d) for an elastic Earth with a fluid core
and an equilibrium ocean. The observed period is about 433\,d. So $\Aw=A_m$ is within 1\% of the
complete elastic value, and the remaining 8\,d is attributed to anelasticity.

%=======================================================================
\section{What the equation does: analytic solutions for testing}
\label{sec:behaviour}
%=======================================================================

The solutions below come from eq.~\eqref{eq:reduced} with $\K=\kappa\Id$ and no damping, or from its complex form
$m=\omega_1+i\omega_2$, in which $\E$ becomes multiplication by $i$. They double as test cases for a
solver.

\paragraph{Constant excitation.} With $\bw_s:=(\Id-\K)^{-1}\bq^\zeta$ the static answer,
$\bw_\perp(t)=\bw_s+\exp[\se\E(\Id-\K)t]\,(\bw_0-\bw_s)$. The rotation vector circles the static
answer at the Chandler rate, prograde, and never settles. A load switched on at $t=0$ from rest
excites a free wobble whose radius equals the static response.

\paragraph{Periodic excitation.} For $\bq^\zeta\propto e^{i\nu t}$ ($\nu>0$ prograde, $\nu<0$ retrograde),
\begin{equation}
  \frac{m_{\rm dynamic}}{m_{\rm static}}=\frac{1}{1-\nu/\scw},
  \label{eq:ratio}
\end{equation}
which is a resonance on the prograde side only (isotropic approximation, using each configuration's $T_{\rm CW}$):
\begin{center}
\small
\begin{tabular}{@{}lcccc@{}}
\toprule
Forcing period & 30\,d & 182.6\,d & 365.25\,d & $\to\infty$\\
\midrule
prograde, $T_{\rm CW}=476$\,d ($\Aw=I_1$) & $-0.067$ & $-0.62$ & $-3.30$ & $1$\\
retrograde, $T_{\rm CW}=476$\,d & $0.059$ & $0.28$ & $0.43$ & $1$\\
prograde, $T_{\rm CW}=422$\,d ($\Aw=A_m$) & $-0.077$ & $-0.76$ & $-6.42$ & $1$\\
retrograde, $T_{\rm CW}=422$\,d & $0.066$ & $0.30$ & $0.46$ & $1$\\
\bottomrule
\end{tabular}
\end{center}
At monthly periods the polar motion is an order of magnitude smaller than static theory predicts
and of opposite sign. At annual periods the prograde part is amplified several-fold, by an amount
sensitive to $T_{\rm CW}$: $-5.4$ at the observed 433\,d. These factors multiply the rotational \emph{correction} to sea level, not
sea level itself.

\paragraph{Linear trend.} For $\bq^\zeta=\dot\bq^\zeta t$ from rest,
\begin{equation}
  \bw_\perp(t)=\bw_s(t)+\frac{1}{\scw}\,\E^T\,\dot\bw_s+\bw_{\rm free}(t),\qquad
  \bw_s(t)=\frac{\dot\bq^\zeta t}{1-\kappa},
  \label{eq:trend}
\end{equation}
where $\bw_{\rm free}$ is a free wobble of the same amplitude as the constant offset. The trend is
exactly the static one. The offset $\E^T\dot\bw_s/\scw$ is \emph{perpendicular} to the trend: the pole
is displaced sideways from the static path by the distance the static pole travels in
$T_{\rm CW}/2\pi$. That is 76\,d for $T_{\rm CW}=476$\,d and 67\,d for 422\,d. It is not a lag along
the path, which would be parallel to the trend. Decadal trends are therefore unaffected, while
sub-annual variability is.

%=======================================================================
\section{Solving the forward problem iteratively}
\label{sec:solve}
%=======================================================================

\subsection{Time discretisation}
\label{sec:timestep}

The elastic theory has no dissipation, so its free wobble never decays. The real Chandler wobble
has $Q\approx100$, mainly from mantle anelasticity. The solver therefore integrates
\begin{equation}
  \dot\bw_\perp=\mathsf{L}\,(\bw_\perp-\bq),\qquad \mathsf{L}=\se\E-\gamma\Id,
  \label{eq:damped}
\end{equation}
where $\gamma\ge0$ is a phenomenological damping of the Munk--MacDonald type. It leaves the static
limit unchanged. With $\K=\kappa\Id$ the free mode is $e^{(i\scw-\gamma(1-\kappa))t}$, so
$Q=\se/2\gamma$. The recommended default for applications is $Q\approx100$. Setting $\gamma=0$
recovers the purely elastic theory, which is used for the checks of \S\ref{sec:checks}. Damping
preserves reciprocity (\S\ref{sec:adjoint}).

Apply the trapezoidal rule on a grid $t_n=n\Delta t$, with $\bw_n=\bw_\perp(t_n)$, $\bq_n=\bq(t_n)$ and
$h=\Delta t/2$:
\begin{equation}
  \boxed{\;\bw_{n+1}=\R\,\bw_n+\Sm\,(\bq_n+\bq_{n+1}),\qquad
  \R=(\Id-h\mathsf{L})^{-1}(\Id+h\mathsf{L}),\qquad \Sm=-(\Id-h\mathsf{L})^{-1}h\mathsf{L}\;}
  \label{eq:step}
\end{equation}
For $\gamma=0$ and $a=\se\Delta t/2$, these are
$\R=[(1-a^2)\Id+2a\E]/(1+a^2)$ and $\Sm=(a^2\Id-a\E)/(1+a^2)$.
Properties:
\begin{itemize}
  \item $\R+2\Sm=\Id$, so a steady state $\bw=\bq$ is preserved exactly. For $\gamma=0$, $\R$ is an
  exact rotation by $2\arctan a$, so the scheme adds no numerical damping of its own.
  \item $\|\Sm\|<1$ at every $\Delta t$, with $\|\Sm\|=a/\sqrt{1+a^2}$ for $\gamma=0$. This follows
  because $\mathsf{L}$ is normal with eigenvalues of non-positive real part. It is what makes the
  inner iteration of \S\ref{sec:schemeA} converge unconditionally.
  \item \textbf{Accuracy.} With feedback, the discrete free mode is a rotation by
  $2\arctan(a_{\rm CW})$ per step, $a_{\rm CW}=\scw\Delta t/2$, so the period is too long by about
  $a_{\rm CW}^2/3$: 0.09\% for weekly steps and 1.7\% for monthly steps (429.3 against 422.2\,d in the default configuration). The
  phase error accumulates over a run. In a 20-year test with weekly steps the maximum error in
  $\bw$ was 3.5\%. For long runs, use $\Delta t$ of a week or less, or use the exact propagator of
  \S\ref{sec:schemeB}.
  \item Backward Euler, $\bw_{n+1}=[(\Id+b\E)\bw_n+(b^2\Id-b\E)\bq_{n+1}]/(1+b^2)$ with
  $b=\se\Delta t$ (for $\gamma=0$), damps the wobble numerically but tends to the static solution
  as $\Delta t\to\infty$. It is useful only as a check against the static solver.
  \item \textbf{Units.} \texttt{pyslfp} is non-dimensional. Compute $\se$ from
  \texttt{rotation\_frequency}, the two moments and $A_f/I_1=0.113$ (PREM; not in \texttt{pyslfp}), and express $\Delta t$ in units of
  \texttt{time\_scale}.
\end{itemize}

\subsection{Scheme 1: march in time with an inner fixed-point iteration (recommended)}
\label{sec:schemeA}

The static solver already iterates between the load, the fields and $\bw$. Only its rotation update
changes. Given the converged $\bw_n$ and $\bq_n$ from the previous step and the load $\zeta_{n+1}$:
\begin{center}
\fbox{\begin{minipage}{0.93\linewidth}
\begin{enumerate}
  \item Initialise $\bw^{(0)}=\bw_n$ and the load $\sigma^{(0)}$ (a warm start from the converged
  $\sigma_n$ is better than a cold one).
  \item For $k=0,1,\dots$ until the load converges:
  \begin{enumerate}
    \item $\bm\psi_{21}=r\,\bw^{(k)}$;
    \item $u_{lm}=h_l\sigma^{(k)}_{lm}$, $\phi_{lm}=k_l\sigma^{(k)}_{lm}$, adding $h^T_2\bm\psi_{21}$
    and $k^T_2\bm\psi_{21}$ at $(2,1)$;
    \item $\bq^{(k)}=\beta\bm\phi_{21}$;
    \item \textbf{the only change:} $\bw^{(k+1)}=\R\bw_n+\Sm(\bq_n+\bq^{(k)})$, in place of the
    static $\bw^{(k+1)}=\bq^{(k)}$;
    \item add $r\bw^{(k+1)}$ to $\bm\phi_{21}$, form $\xi$ from eq.~\eqref{eq:sl}, and set
    $\sigma^{(k+1)}=\zeta_{n+1}+\rho_wC\xi$.
  \end{enumerate}
  \item Store $\bw_{n+1}=\bw^{(k+1)}$ and $\bq_{n+1}=\bq^{(k)}$, the latter computed with $\phi$ excluding $\psi$.
\end{enumerate}
\end{minipage}}
\end{center}
In \texttt{pyslfp} this is a copy of \texttt{solve\_generalised\_equation} with the rotation
update replaced:
\begin{verbatim}
    q = i * potential_change_lm.coeffs[:, 2, 1]
    if angular_momentum_change is not None:
        q -= m * angular_momentum_change
    angular_velocity_change = R @ w_old + S @ (q_old + q)
\end{verbatim}
\sloppy The copy takes \texttt{w\_old} and \texttt{q\_old} as inputs and returns the converged \texttt{q}. The \texttt{angular\_momentum\_change} lines carry the generalised forcing $\bk$ of \S\ref{sec:gen}
and are inactive in a forward run.
The static solver is the case $\R=0$, $\Sm=\Id$ with the $\bq_n$ term removed.

\paragraph{Convergence.} With the ocean converged, the rotational part of the inner iteration is
$\bw\mapsto\R\bw_n+\Sm(\bq_n+\bq^\zeta_{n+1}+\K\bw)$. It contracts by
$\|\Sm\K\|\le\|\Sm\|\,\kappa_{\max}<0.37$ at every step size and forcing frequency, and more
strongly for small steps. In a test with the $\K$ of eq.~\eqref{eq:Knum} and weekly steps, the
iteration reached $10^{-14}$ in about 10 iterations. In practice the ocean-load iteration, which is
the same as in the static solver, dominates the iteration count.

\paragraph{Why this is the recommended solver.} It is the smallest change to working code. It
never iterates over time: the problem is causal, so step $n+1$ needs only the two 2-vectors
$\bw_n,\bq_n$. It is also the scheme that survives the extensions of \S\ref{sec:ext}: moving
shorelines (a nonlinear static solve per step), damping or anelastic memory (extra state variables
carried alongside $\bw_n$), and $\bh$ or $\bGam$ (additional terms in $\bq$).

\subsection{Scheme 2: exact elimination of the rotation (verification and fast path)}
\label{sec:schemeB}

Because map (i) of \S\ref{sec:problem} is linear and time-invariant, $\bw$ can be eliminated exactly.
This needs no change to \texttt{pyslfp}.
\begin{enumerate}
  \item \textbf{Rotational basis (two solves).} For $\mu=1,2$, run the \emph{static} generalised
  solver with no load and rotational forcing $\bk=\be_\mu$ (\S\ref{sec:gen}), used here purely as a
  numerical device. Since $\bk$ enters only the rotation row, the resulting
  fields $X[\be_\mu]=(\xi,\bu,\phi)$ are those of the rotation
  $\bw^{(\mu)}$ that the solver returns. With $\W=[\bw^{(1)}\ \bw^{(2)}]$, the static relation
  $\bw=\K\bw-\bk/(I_3-I_1)$ gives
  \begin{equation}
    \K=\Id+\frac{\W^{-1}}{I_3-I_1},\qquad
    X^\alpha=\sum_\mu X[\be_\mu]\,(\W^{-1})_{\mu\alpha},
    \label{eq:basis}
  \end{equation}
  where $X^\alpha$ is the full response to a unit prescribed rotation $\bw_\perp=\be_\alpha$ with the ocean
  responding. This is how eq.~\eqref{eq:Knum} was computed.
  \item \textbf{Non-rotating solves.} For each $t_n$, solve with $\zeta_n$ and
  \texttt{rotational\_feedbacks=False}. This gives $X^\zeta_n$ and $\bq^\zeta_n=\beta\bm\phi^\zeta_{21}$.
  These solves are independent and can be run in parallel.
  \item \textbf{Rotation.} Integrate eq.~\eqref{eq:reduced}. With the trapezoidal rule the step is
  the $2\times2$ solve $(\Id-\Sm\K)\bw_{n+1}=\R\bw_n+\Sm(\bq_n+\bq^\zeta_{n+1})$, where
  $\bq_n=\bq^\zeta_n+\K\bw_n$. This reproduces Scheme~1 exactly: they agreed to $5\times10^{-14}$ in
  the test. Alternatively, with $\mathsf{M}=\mathsf{L}(\Id-\K)$, $\mathbf{f}=-\mathsf{L}\bq^\zeta$ taken
  piecewise linear, $\mathsf{P}_0=e^{\mathsf{M}\Delta t}$, $\mathsf{P}_1=\mathsf{M}^{-1}(\mathsf{P}_0-\Id)$ and
  $\mathsf{P}_2=\mathsf{M}^{-1}(\mathsf{P}_1/\Delta t-\Id)$, the step
  $\bw_{n+1}=\mathsf{P}_0\bw_n+\mathsf{P}_1\mathbf{f}_n+\mathsf{P}_2(\mathbf{f}_{n+1}-\mathbf{f}_n)$
  is exact and has no period error.
  \item \textbf{Reassemble:} $X_n=X^\zeta_n+\sum_\alpha\omega_{\alpha,n}X^\alpha$.
\end{enumerate}
The cost is two extra solves, and changing the damping, the time step or the initial wobble
requires no further sea-level solves. It fails as soon as the problem is nonlinear (moving
shorelines) or the response has memory in (i) (viscoelasticity). That is why it is the check
rather than the main solver.

\subsection{Why not iterate over the whole time history?}
\label{sec:waveform}

The alternative (Method~B of the earlier companion note) guesses $\xi(\bx,t)$ for all $t$, builds the
excitation history, integrates the wobble ODE over $[0,T]$, recomputes $\xi$ everywhere, and
repeats. Its fixed point is the same, but it has two problems.

\paragraph{It diverges in the band of interest.} If the propagator uses $\se$, the loop must
generate the whole feedback $\K$. One lap then has gain $\kappa/|1-\nu/\se|$ at frequency
$\nu$. This exceeds one for periods between $T_e/(1+\kappa)=206$\,d and $T_e/(1-\kappa)=394$\,d ($\Aw=A_m$),
which contains the annual band. On a finite window the discrete iteration matrix is block lower
triangular with spectral radius $\|\Sm\K\|<1$, so it converges formally. Its non-normal transient
growth, of order $(\kappa\se T)^k/k!$, is fatal in practice. In a 20-year test with annual forcing
the iterates grew by $10^{20}$ in 35 iterations.

\paragraph{The Chandler-propagator variant is easy to get wrong.} Folding the solid-Earth feedback
into the propagator, $\dot\bw=\scw^{\rm solid}\E(\bw-\bchi')$, is legitimate only with
$\bchi'=\bq^\zeta_{\rm load}/(1-\kappa)$. Here $\bq^\zeta_{\rm load}$ excludes the tidal term
$k^T_2\bm\psi_{21}$, although that term must still be included in the \emph{fields} $\bu,\phi$.
Using the total $\bq$ counts the solid-Earth rotational deformation twice, because
$\scw\E((1-\kappa)\bw-\bq^\zeta)\neq\se\E((1-\kappa)\bw-\bq^\zeta)$. Even when done correctly, the
residual ocean gain of about $0.05/|1-\nu/\scw|$ still diverges within about 5\% of the Chandler
period.

\paragraph{Conclusion.} For this linear and causal problem, with its instantaneous response, marching
(Scheme~1) is better in every respect. Nonlinearity does not change this, because the problem stays
causal and Scheme~1 simply solves a nonlinear static problem at each step.

\subsection{Verification checklist}
\label{sec:checks}

\begin{enumerate}
  \item \textbf{Static limit.} For slowly varying $\zeta$, Scheme~1 tends to the static
  \texttt{pyslfp} solution. Backward Euler with very large $\Delta t$ reproduces it to solver
  tolerance.
  \item \textbf{Free wobble.} With $\gamma=0$, $\zeta=0$ and $\bw_0\neq\mathbf{0}$, the motion is a
  closed ellipse with period $2\pi/(\se\sqrt{\det(\Id-\K)})$, as tabulated in \S\ref{sec:spectral},
  with no growth or decay. With $\gamma>0$ it decays with the chosen $Q$.
  \item \textbf{$\K$.} $\K=\K^T$ and $\mathrm{tr}\,\E(\Id-\K)=0$ to round-off (both confirmed at
  $l_{\max}=64$). With $C=0$, $\K=\kappa\Id$.
  \item \textbf{Scheme~1 = Scheme~2} for the same $\Delta t$, to solver tolerance.
  \item \textbf{Analytic cases.} The ratio of eq.~\eqref{eq:ratio} under periodic forcing, and the
  perpendicular offset of eq.~\eqref{eq:trend} under a trend.
  \item \textbf{Order of accuracy.} Second-order convergence in $\Delta t$.
\end{enumerate}

%=======================================================================
\section{Extensions: where each enters}
\label{sec:ext}
%=======================================================================

\paragraph{Relative angular momentum and torques.} These enter only through the excitation of
eq.~\eqref{eq:chi}: replace $\bq$ by $\bchi$. At leading order this means $\bq\to\bq+\bh_\perp/(I_3-I_1)$
and $\omega_3=-(j_3+h_3)/I_3$. This is AA24's $\bk$ with $\bk=-\bh$, which \texttt{pyslfp} already
carries as \texttt{angular\_momentum\_change}. A torque enters as
$\E\bGam/[\Omega(I_3-I_1)]$, rotated by $90^\circ$. The terms $\dot\bh$, and $\dot\bj$ if wanted,
are $O(\epsilon)$ and should be added together or not at all. If they are added, the cleanest state
variable is the equatorial angular momentum $\bl=I_1\bw+\bj+\bh$, for which eq.~\eqref{eq:eqE} is
exactly
\begin{equation}
  \dot\bl=\Omega\E\,(I_3\bw-\bl)+\bGam,\qquad \bw=(\bl-\bj-\bh)/I_1 .
  \label{eq:ell}
\end{equation}
Scheme~1 carries over directly, marching $\bl$ instead of $\bw$, with $\bh$ including
$\bh_{\rm core}$ if the core is treated this way rather than through $\Aw$. The inner update
$\bw\leftarrow(\bl-\bj-\bh)/I_1$ contracts by about $\delta\kappa\sim10^{-3}$. Retaining these
$O(\epsilon)$ terms is only worthwhile once the larger neglected effects of \S\ref{sec:order} have
been modelled.

\paragraph{Anelasticity.} The phenomenological damping of eq.~\eqref{eq:damped} is built into the
time step. Do not also use an observed or Chandler frequency in the propagator
(\S\ref{sec:waveform}). The physical alternative is anelastic, frequency-dependent Love numbers.
These give map (i) memory, which in the time domain means internal (relaxation) variables marched
alongside $\bw_n$. They would also supply the remaining $\approx8$\,d of the Chandler period.
Scheme~1 accommodates them, but Scheme~2 does not. Anelastic Love numbers are expected to preserve
reciprocity, being a dissipative linear system, but this has not been checked.

\paragraph{Fluid core, beyond leading order.} The treatment through $\Aw$ (\S\ref{sec:liouville})
captures the core's main effect. A fuller treatment includes wobble-specific Love numbers and
core--mantle coupling (Smith \& Dahlen 1981), and ocean and atmospheric transfer functions (Wahr 1982).

\paragraph{Triaxiality.} The forward solver needs only eq.~\eqref{eq:components} with $I_1\neq I_2$,
which gives a $2\times2$ ODE with different coefficients. The analogues of $\R$ and $\Sm$ follow
from the same trapezoidal rule. Reciprocity still holds; the adjoint has the rotation reversed and $I_1$, $I_2$
exchanged in the gyroscopic term (\S\ref{sec:adjoint}).

\paragraph{Observational conventions.} See \S\ref{sec:data}. These affect the observation
operators, not the forward solver.

%=======================================================================
\section{The generalised problem and reciprocity}
\label{sec:recip}
%=======================================================================

This section extends AA24 \S\S2.7 and 3.2 to the dynamic problem. The aim is that the forward
solver of \S\ref{sec:solve}, run backwards in time, is also the adjoint solver.

\subsection{The generalised dynamic problem}
\label{sec:gen}

Observables other than sea level need forcings conjugate to each unknown. AA24 add three:
\begin{itemize}
  \item a surface force density $\bt$, paired with $\bu$;
  \item a fictitious surface mass $\zeta_\phi$, paired with $\phi$, which sources the potential but
  exerts no force;
  \item a rotational forcing $\bk$, paired with $\bw$.
\end{itemize}
A real load is the combination $(\bu'\cdot\nabla\Phi+\phi')\sigma$. All forcings may now depend on
time. The generalised dynamic problem is
\begin{equation}
  \boxed{\;\begin{aligned}
  &A(\bu,\phi\,|\,\bu',\phi')+\int_{\pM}(\bu'\cdot\nabla\Phi+\phi')\,\sigma\dd S
  +\int_{\pM}(\bt\cdot\bu'+\zeta_\phi\,\phi')\dd S\\
  &\qquad+\int_M\rho\,\bu'\cdot\nabla\psi\dd V+(\It\bw+\bN\dot\bw-\bj+\bk)\cdot\bw'=0
  \end{aligned}\;}
  \label{eq:gen}
\end{equation}
with $\bw(0)=\bw_0$. Here $\sigma$, $\xi$ and $\int\sigma\dd S=0$ are as in eq.~\eqref{eq:sl}. Because
$\zeta_\phi$ counts as mass for the potential, it also counts in the inertia (AA24 eq.~56):
\begin{equation}
  \bw'\cdot\bj+\int_M\rho\,\bu\cdot\nabla\psi'\dd V+\int_{\pM}(\sigma+\zeta_\phi)\,\psi'\dd S=0 .
  \label{eq:inertiagen}
\end{equation}
\paragraph{The matrix $\bN$.} $\bN$ is any matrix with no axial part. This covers three cases. The
undamped case with $A_{w,1}=A_{w,2}$ has equatorial block $(\Aw/\Omega)\E$. The triaxial case is
eq.~\eqref{eq:N}. The phenomenological damping of eq.~\eqref{eq:damped} has equatorial block
$(I_3-I_1)(\se\E+\gamma\Id)/(\se^2+\gamma^2)$, which reproduces $\dot\bw=(\se\E-\gamma\Id)(\bw-\bq)$.
In every case the antisymmetric part of $\bN$ is gyroscopic and its symmetric part is dissipative.

\paragraph{Correspondence with \texttt{pyslfp}.} \sloppy The code supports $\bt=\zeta_u\nabla\Phi$
(\texttt{displacement\_load}), $\zeta_\phi$ (\texttt{gravitational\_potential\_load}) and $\bk$
(\texttt{angular\_momentum\_change}). Their responses enter through the generalised Love numbers of
AA24 eq.~62. The rotational row is unchanged, with $\bq=\beta\bm\phi_{21}-\bk_\perp/(I_3-I_1)$, where
$\bm\phi_{21}$ now includes the potential sourced by $\zeta_\phi$. Scheme~1 therefore solves
eq.~\eqref{eq:gen} with the same one-line change as before.

\paragraph{Solvability.} For a spherically symmetric reference, the body force $-\rho\nabla\psi$
exerts no net force or torque, so the condition is AA24's eq.~30 at each $t$:
$\int_{\pM}(\bt-\zeta_\phi\nabla\Phi)\cdot(\ba+\bb\times\bx)\dd S=0$. Equivalently, the observable
must be invariant under rigid-body motions.

\subsection{The reciprocity theorem}
\label{sec:theorem}

Let $(\bu,\phi,\bw)$ solve eq.~\eqref{eq:gen} with loads $(\zeta,\bt,\zeta_\phi,\bk)$ and
$\bw(0)=\bw_0$. Let $(\bu^\dagger,\phi^\dagger,\bw^\dagger)$ solve it with loads
$(\zeta^\dagger,\bt^\dagger,\zeta^\dagger_\phi,\bk^\dagger)$, with $\bN$ replaced by
\begin{equation}
  \bN^\dagger:=-\bN^T,
  \label{eq:Ndag}
\end{equation}
and with the \emph{terminal} condition $\bw^\dagger(T)=\mathbf{0}$. Four steps give the theorem.
\begin{enumerate}
  \item \textbf{Instantaneous identity.} Test the first problem with the dagger fields and the
  second with the undaggered fields, and subtract. The $A$ terms cancel by symmetry, and so do the
  $\It$ terms. Eliminating $\bj$ and $\bj^\dagger$ with eq.~\eqref{eq:inertiagen} makes the volume
  integrals cancel, leaving at each $t$
  \begin{equation}
    \int_{\pM}(\bu^\dagger\!\cdot\!\nabla\Phi+\phi^\dagger+\psi^\dagger)\sigma\dd S+X
    -\int_{\pM}(\bu\!\cdot\!\nabla\Phi+\phi+\psi)\sigma^\dagger\dd S-X^\dagger
    =-\big(\bw^\dagger\!\cdot\bN\dot\bw-\bw\cdot\bN^\dagger\dot\bw^\dagger\big),
    \label{eq:inst}
  \end{equation}
  where $X=\int_{\pM}[\bt\cdot\bu^\dagger+\zeta_\phi(\phi^\dagger+\psi^\dagger)]\dd S+\bk\cdot\bw^\dagger$,
  and $X^\dagger$ is the same with the roles exchanged.
  \item \textbf{Total derivative.} By eq.~\eqref{eq:Ndag},
  $\bw\cdot\bN^\dagger\dot\bw^\dagger=-\dot\bw^\dagger\cdot\bN\bw$. The right-hand side of
  eq.~\eqref{eq:inst} is therefore $-\frac{\dd}{\dd t}(\bw^\dagger\cdot\bN\bw)$, for any $\bN$.
  This is why the adjoint carries $-\bN^T$. When $\bN$ is antisymmetric, $\bN^\dagger=\bN$.
  \item \textbf{Integrate over $[0,T]$.} The boundary term is
  $-[\bw^\dagger\cdot\bN\bw]_0^T=\bw^\dagger(0)\cdot\bN\bw_0$, since $\bw^\dagger(T)=\mathbf{0}$.
  \item \textbf{Convert to sea level.} Use eq.~\eqref{eq:sl}. The constants $\Phi_g,\Phi_g^\dagger$
  vanish against $\int\sigma^\dagger=\int\sigma=0$, and the ocean terms
  $\rho_w\int C\xi\xi^\dagger$ cancel. Finally divide by $-g$.
\end{enumerate}
The result is
\begin{multline}
  \int_0^T\!\Big\{\int_{\pM}\xi^\dagger\zeta\dd S-\frac1g\int_{\pM}\big[\bt\cdot\bu^\dagger+\zeta_\phi(\phi^\dagger+\psi^\dagger)\big]\dd S-\frac1g\,\bk\cdot\bw^\dagger\Big\}\dd t\\
  =\int_0^T\!\Big\{\int_{\pM}\xi\,\zeta^\dagger\dd S-\frac1g\int_{\pM}\big[\bt^\dagger\cdot\bu+\zeta^\dagger_\phi(\phi+\psi)\big]\dd S-\frac1g\,\bk^\dagger\cdot\bw\Big\}\dd t
  -\frac1g\,\bw^\dagger(0)\cdot\bN\bw_0 .
  \label{eq:recipgen}
\end{multline}
This is AA24's eq.~57, integrated over time, plus an initial-state term. With only direct loads and
$\bw_0=\mathbf{0}$ it reduces to
\begin{equation}
  \boxed{\;\int_0^T\!\!\int_{\pM}\xi^\dagger\,\zeta\dd S\dd t=\int_0^T\!\!\int_{\pM}\xi\,\zeta^\dagger\dd S\dd t\;}
  \label{eq:recip}
\end{equation}
Since $\psi=-[\bx\times(\bO\times\bx)]\cdot\bw$, the $\zeta_\phi\psi$ terms can be moved into the
rotational forcing by replacing $\bk$ with $\bk-\int_{\pM}\zeta_\phi\,\bx\times(\bO\times\bx)\dd S$
(AA24 eq.~58). This leaves $\zeta_\phi$ paired with $\phi$ alone.

\subsection{The adjoint problem in practice}
\label{sec:adjoint}

Write $\tau=T-t$ and $\tilde\bw(\tau)=\bw^\dagger(T-\tau)$, and similarly for the other fields and
loads. Then $\bN^\dagger\dot\bw^\dagger=-\bN^T\dot\bw^\dagger=\bN^T\dd\tilde\bw/\dd\tau$, and the
adjoint problem becomes
\begin{equation}
  \text{eq.~\eqref{eq:gen} in }\tau,\ \text{with }\bN\to\bN^T,\ \text{loads }(\tilde\zeta,\tilde\bt,\tilde\zeta_\phi,\tilde\bk),\ \tilde\bw(0)=\mathbf{0}.
  \label{eq:adjrev}
\end{equation}
In words: \emph{run the forward generalised solver forwards in $\tau$, on time-reversed adjoint loads,
with $\bN$ replaced by its transpose, then reverse the output.}
\begin{center}
\small
\begin{tabular}{@{}lll@{}}
\toprule
Case & $\bN^T$ in terms of $\bN$ & Adjoint solver\\
\midrule
$A_{w,1}=A_{w,2}$, undamped & $-\bN$ & rotation reversed ($\se\to-\se$)\\
damped & gyroscopic part reversed, damping kept & rotation reversed, same $\gamma$\\
triaxial & $-\bN$ with $A_{w,1}\leftrightarrow A_{w,2}$ & reversed, $A_{w,1},A_{w,2}$ swapped\\
\bottomrule
\end{tabular}
\end{center}
With damping, the adjoint is therefore as stable in $\tau$ as the forward problem is in $t$. The
static parts ($\It$, $\K$, the Love numbers and the ocean) are self-adjoint and unchanged.

\paragraph{In Scheme~1.} The step matrices $\R$ and $\Sm$ are rational functions of the matrix
$\mathsf{L}=\se\E-\gamma\Id$. Replacing $\mathsf{L}$ by $\mathsf{L}^T$ replaces them by $\R^T$ and
$\Sm^T$, which for $\gamma=0$ is $a\to-a$. The adjoint solver is Scheme~1 with $\R^T,\Sm^T$, run on
reversed loads.

\paragraph{Green's function.} Let $G_{\bN}(\bx_2,\bx_1;s)$ be the sea level at $\bx_2$ a time $s$ after
a unit impulsive load at $\bx_1$. Then eq.~\eqref{eq:recip} gives
\begin{equation}
  G_{\bN}(\bx_2,\bx_1;s)=G_{\bN^T}(\bx_1,\bx_2;s).
  \label{eq:green}
\end{equation}
Swapping source and receiver is compensated by reversing the rotation, which is Onsager--Casimir
reciprocity. Integrating over $s$ recovers AA24's static symmetry, because the static response does
not depend on $\bN$.

\subsection{Sensitivity kernels and adjoint loads}
\label{sec:kernels}

Suppose a functional $J$ has first-order perturbation
\begin{equation}
  \delta J=\int_0^T\!\Big\{\int_{\pM}\zeta^\dagger\delta\xi\dd S-\frac1g\int_{\pM}\big[\bt^\dagger\cdot\delta\bu+\zeta_\phi^\dagger(\delta\phi+\delta\psi)\big]\dd S-\frac1g\,\bk^\dagger\cdot\delta\bw\Big\}\dd t .
  \label{eq:dJ}
\end{equation}
Apply eq.~\eqref{eq:recipgen} with a forward problem driven by $\delta\zeta$ and $\delta\bw_0$ only. This gives
\begin{equation}
  \delta J=\int_0^T\!\!\int_{\pM}\xi^\dagger\,\delta\zeta\dd S\dd t+\frac1g\,\bw^\dagger(0)\cdot\bN\,\delta\bw_0 .
  \label{eq:kernel}
\end{equation}
The ice kernel is $K(\bx,t)=\rho_i(1-C)\,\xi^\dagger(\bx,t)$. The second term is the kernel for the initial
wobble, which an inversion must treat as a nuisance parameter unless the run starts from rest
(\S\ref{sec:problem}). The adjoint loads for common observables follow AA24 \S4 and the altimetry
row of Note~1, each multiplied by a time weight $w(t)$ ($\delta(t-t_o)$ at an instant, a boxcar for
a monthly mean):
\begin{center}
\small
\begin{tabular}{@{}lllll@{}}
\toprule
Observable (times $w(t)$) & $\zeta^\dagger$ & $\bt^\dagger$ & $\zeta^\dagger_\phi$ & $\bk^\dagger$\\
\midrule
$\xi(\bx_0)$ & $\delta(\bx,\bx_0)$ & $\mathbf{0}$ & $0$ & $\mathbf{0}$\\
$\bn\cdot\bu(\bx_0)$ & $0$ & $-g\,\delta(\bx,\bx_0)\bn$ & $0$ & $\mathbf{0}$\\
$\phi_{lm}$ (AA24 eq.~66) & $0$ & $\mathbf{0}$ & $-g\,Y_{lm}$ & $\int_{\pM}\zeta^\dagger_\phi\,\bx\times(\bO\times\bx)\dd S$\\
$\int_{\pM}v\,\Delta\mathrm{SSH}^*\dd S$ & $v$ & $-g\,v\,\bn$ & $0$ & $\int_{\pM}v\,\bx\times(\bO\times\bx)\dd S$\\
$\be_\alpha\cdot\bw$ (polar motion) & $0$ & $\mathbf{0}$ & $0$ & $-g\,\be_\alpha$\\
\bottomrule
\end{tabular}
\end{center}
In the fourth row, $v$ is the spatial averaging weight and
$\Delta\mathrm{SSH}^*=\xi+\bn\cdot\bu+\psi/g$. This is the geocentric sea-surface height with only the
direct centrifugal term removed (Lickley et al.\ 2018). It is a modelling convention, not the operational
pole-tide correction. For products corrected in the standard way, eq.~\eqref{eq:dcorr} applies instead.
The operator is then the non-rotating one, whose adjoint has $\zeta^\dagger=v$ and
$\bt^\dagger=-gv\bn$, solved with rotational feedback switched off (\texttt{rotational\_feedbacks=False}).
Setting $\bk^\dagger=\mathbf{0}$ in the rotating adjoint is not enough, since that retains the
load-driven rotation that cancels from the data.
The $\bk^\dagger$ entries in rows three and four are non-zero only for $(l,m)=(2,\pm1)$, plus $l=0$
and $(2,0)$ if $\omega_3$ is kept.

\paragraph{Causality.} If $w(t)=0$ for $t>t_o$, the adjoint fields vanish for $t>t_o$: in $\tau$ nothing
forces them until the observation time is reached. The kernel therefore automatically ignores loads
applied after the measurement. Before $t_o$ the kernel contains a free wobble at the Chandler
period, decaying into the past with an e-folding time of about 40\,yr for $Q\approx100$, and not
at all if $\gamma=0$ (\S\ref{sec:worries}).

\subsection{Discrete reciprocity}
\label{sec:discrete}

The trapezoidal scheme satisfies reciprocity \emph{exactly}, with no separately derived discrete
adjoint, provided time integrals are evaluated by pairing interval means:
\begin{equation}
  \int_0^T f\,g\dd t\;\to\;\sum_n\Delta t\,\bar f_{n+\frac12}\,\bar g_{n+\frac12},\qquad
  \bar f_{n+\frac12}=\tfrac12(f_n+f_{n+1}).
  \label{eq:pairing}
\end{equation}
The reason is that the trapezoidal step is exactly eq.~\eqref{eq:gen} applied to interval means,
with $\dot\bw$ replaced by $(\bw_{n+1}-\bw_n)/\Delta t$. The instantaneous relations are linear, so
the means satisfy them exactly. The gyroscopic terms telescope for any $\bN$:
\begin{equation*}
  \bar\bw^\dagger\cdot\bN(\bw_{n+1}-\bw_n)-\bar\bw\cdot\bN^\dagger(\bw^\dagger_{n+1}-\bw^\dagger_n)
  =\bw^\dagger_{n+1}\cdot\bN\bw_{n+1}-\bw^\dagger_n\cdot\bN\bw_n .
\end{equation*}
This was confirmed on the reduced coupled model (the $\K$ of eq.~\eqref{eq:Knum}, a random symmetric
non-rotating response, and the rotational coupling implied by \S\ref{sec:K}). The two sides of
eq.~\eqref{eq:recip} agreed to $10^{-14}$, with and without damping, at every $\Delta t$. Pairing
nodal values with trapezoidal weights instead agrees only to $O(\Delta t^2)$. In practice, define
observables as interval means, which is natural for monthly GRACE and altimetry products. Adjoints
from the same solver then pass dot-product tests at machine precision.

\subsection{Things to be aware of}
\label{sec:worries}

\begin{enumerate}
  \item \textbf{Untested in the full code.} Every numerical test used the reduced model built from a
  real $\K$. Scheme~1 has not yet been implemented in \texttt{pyslfp}. Its combined ocean-plus-rotation
  inner iteration is argued to converge but has not been demonstrated.
  \item \textbf{Near-resonant results depend on the Chandler period.} With $\Aw=A_m$ the model gives
  422\,d, against 425.5\,d for a complete elastic calculation and about 433\,d observed. The annual
  prograde factor of eq.~\eqref{eq:ratio} is nevertheless $-6.4$ at 422\,d and $-5.4$ at 433\,d.
  Annual-band rotational corrections therefore carry an uncertainty of order 20\% unless the
  remaining anelastic lengthening is modelled. Trends are unaffected.
  \item \textbf{Long memory.} The elastic free wobble never decays. Any excitation (a load switched
  on, an inconsistent initial state) then persists for ever, and adjoint kernels extend over the whole
  past at full amplitude. This is why $Q\approx100$ damping is the recommended default
  (\S\ref{sec:timestep}). Even then the e-folding time is about $QT_{\rm CW}/\pi\approx40$\,yr, so
  $\bw_0$ and loads before the window genuinely matter.
  \item \textbf{Real polar motion is not load-driven.} The observed wobble is excited mainly by the
  atmosphere and ocean, and its pole tide (a few mm) far exceeds the load-induced one. For corrected
  data this is removed by processing (\S\ref{sec:data}). For uncorrected data it is signal the model
  does not contain.
\end{enumerate}

%=======================================================================
\appendix
\section{Audit of the earlier notes}
\label{app:audit}
%=======================================================================

\paragraph{Checked and correct.} Restatement of AA24 \S\S2--3 (including the Euler-equation
algebra, eqs.~37--49, and the MacCullagh factors). The Liouville equation D3.2 and its equatorial
form D3.3, re-derived and checked symbolically here. The gyroscopic matrix and its antisymmetry
condition. The size estimates in the ordering argument. The axial result. DP1--DP2, DP5 (444\,d) and
$k_s=0.951$, $\kappa=0.314$ (recomputed from the \texttt{pyslfp} constants). The periodic-response
table. The magnitude of the trend offset (70.7\,d). The trapezoidal and backward-Euler matrices, $\R+2\Sm=\Id$,
the Scheme~A contraction bound, and the 450.6\,d discrete period (450.8\,d here). The time-integrated
reciprocity, terminal condition and time reversal. The consistency requirements 1 and 2 of the
companion note's Method~B.

\paragraph{Issues.}
\begin{enumerate}
  \item \label{it:methodB} \textbf{Method~B double counts the solid-Earth feedback (companion note,
  Steps 3--4 and requirement 3).} The ``recommended'' variant pairs a no-ocean Chandler propagator with
  an excitation read from the \emph{total} potential, which already contains $k^T_2\bm\psi_{21}$. It also
  omits the factor $1/(1-\kappa)$. The correct pairing is given in \S\ref{sec:waveform}. Since whole-history
  iteration is not recommended anyway, this affects no conclusion.
  \item \label{it:jdot} \textbf{The $\dot\bj$ refinement (Note~1 \S3.3; companion \S2.7)} is
  algebraically right and harmless, but not systematically better. It is an $O(\epsilon)$ correction to
  the excitation (3\% monthly, 0.3\% annual), of the same order as the neglected relative angular
  momentum of the transported meltwater, and far smaller than the fluid-core and anelastic effects.
  Standard excitation functions omit it. The reciprocity discussion built on it (D3.11--D3.12) is then
  unnecessary. The step-load jump it predicts ($-0.23\%$ of the static response) is real but negligible.
  \item \textbf{The trend offset is not a lag} (Note~1 DP7 and its gloss). The offset is perpendicular
  to the trend, a sideways displacement of the pole path, not a delay along it (\S\ref{sec:behaviour}).
  \item \textbf{The fluid core} is treated in both notes as a later refinement, with the suggestion of
  tuning $\se$. Its relative angular momentum is small, but its gyroscopic term is about 11\% of the
  retained inertial term, the largest effect outside the elastic theory. It is included here through $\Aw=A_m$
  (\S\ref{sec:liouville}), which brings the model period within 1\% of a complete elastic calculation.
  \item \textbf{Pole-tide conventions.} Both notes treat observational conventions as a side issue,
  and Note~1's altimetry row uses the direct-term convention $\Delta\mathrm{SSH}^*$. For data corrected
  in the standard way, the load-driven rotation cancels from the observables (\S\ref{sec:data}).
  \item \textbf{Damping and adjoints (Note~1, ``Physics not included'').} The statement that the adjoint
  ``carries damping of the opposite sign'' is misleading. In reversed time the adjoint is equally
  damped, and rotating-system reciprocity survives (\S\ref{sec:recip}).
  \item \textbf{Triaxiality and discrete adjoints (Note~1).} Note~1 said reciprocity was not
  established for $I_1\neq I_2$, and recommended transposing the time-step recursion to get an exact
  discrete adjoint. Both are superseded. With $\bN^\dagger=-\bN^T$, reciprocity holds for any $\bN$
  (\S\ref{sec:theorem}). With interval-mean pairing, the trapezoidal solver is its own exact discrete
  adjoint (\S\ref{sec:discrete}).
  \item \textbf{Symmetry of $\K$} was conjectured in Note~1 but not proved. It is proved in
  \S\ref{sec:K} and confirmed numerically. The ocean lengthens $T_{\rm CW}$ from 444 to 476\,d.
  \item \textbf{Divergence of the global iteration} (companion \S3.3). The frequency-domain
  argument is a heuristic. On a finite window the iteration converges formally but with
  catastrophic transient growth (\S\ref{sec:waveform}). The conclusion is unchanged.
  \item \textbf{Scheme~B} in both notes asked for a prescribed-$\bw$ solve, which \texttt{pyslfp}
  lacks. Eq.~\eqref{eq:basis} obtains the same basis from the existing angular-momentum forcing.
  \item \textbf{Phase drift.} The notes quote the per-cycle period error of the trapezoidal rule but
  not its accumulation over long runs (\S\ref{sec:timestep}).
  \item \textbf{Sign conventions.} The companion note's eq.~B4 uses the opposite sign for $\phi$ and
  $\psi$. Everything here follows AA24.
  \item \textbf{Constants.} \texttt{pyslfp} uses $\Omega=2\pi/86400$\,s$^{-1}$ (solar day), not the
  sidereal $7.2921\times10^{-5}$\,s$^{-1}$. This changes $\se$ by 0.27\% and $k_s$, $\kappa$ and $\K$ by
0.55\%, which is small but should be corrected in the code.
  \item \textbf{Notation.} Note~1 used $\mathbf{i}$ for the inertia perturbation, $\bGam$ for the
  gyroscopic matrix, $\Delta\mathrm{SL}$ and $\mathcal C$. These are now $\bc$, $\bN$, $\xi$ and $C$, and
  $\bGam$ is the external torque.
\end{enumerate}

\paragraph{Omitted as unnecessary for now.} The demonstration design of the companion note (\S4)
and its $\dot\bj$-based refinements. Neither is wrong as a plan, and the demonstration can be revisited.

\begin{thebibliography}{99}
\bibitem{aa24} Al-Attar, D., Syvret, F., Crawford, O., Mitrovica, J.X.\ \& Lloyd, A.J., 2024.
Reciprocity and sensitivity kernels for sea level fingerprints, \emph{Geophys.\ J.\ Int.}, 236, 362--378.
\bibitem{barnes83} Barnes, R.T.H., Hide, R., White, A.A.\ \& Wilson, C.A., 1983. Atmospheric angular
momentum fluctuations, length-of-day changes and polar motion, \emph{Proc.\ R.\ Soc.\ Lond.\ A}, 387, 31--73.
\bibitem{dahlen76} Dahlen, F.A., 1976. The passive influence of the oceans upon the rotation of the
Earth, \emph{Geophys.\ J.\ R.\ astr.\ Soc.}, 46, 363--406.
\bibitem{desai02} Desai, S.D., 2002. Observing the pole tide with satellite altimetry, \emph{J.\ Geophys.\
Res.}, 107(C11), 3186.
\bibitem{desai15} Desai, S., Wahr, J.\ \& Beckley, B., 2015. Revisiting the pole tide for and from satellite
altimetry, \emph{J.\ Geod.}, 89, 1233--1243.
\bibitem{gross15} Gross, R.S., 2015. Earth rotation variations --- long period, in \emph{Treatise on
Geophysics}, 2nd edn, vol.~3, ed.\ Schubert, G., Elsevier.
\bibitem{heathcote26} Heathcote, D.A., Holland, T., Mag, A.M., Tamisiea, M.E., Coulson, S., Dangendorf, S.,
Lloyd, A.J., Mashayek, A., Mitrovica, J.X.\ \& Al-Attar, D., 2026. A scalable Bayesian framework for modern
sea-level inference, submitted to \emph{Geophys.\ J.\ Int.}
\bibitem{lickley18} Lickley, M.J., Hay, C.C., Tamisiea, M.E.\ \& Mitrovica, J.X., 2018. Bias in estimates of
global mean sea level change inferred from satellite altimetry, \emph{J.\ Climate}, 31, 5263--5271.
\bibitem{munk} Munk, W.H.\ \& MacDonald, G.J.F., 1960. \emph{The Rotation of the Earth}, Cambridge
University Press.
\bibitem{note1} Working notes, September 2026. Note~1 is \texttt{rotation\_alattar\_notation.pdf};
the companion note is \texttt{rotation\_forward\_problem.pdf}.
\bibitem{petit10} Petit, G.\ \& Luzum, B.\ (eds), 2010. \emph{IERS Conventions (2010)}, IERS Technical Note 36,
Verlag des Bundesamts f\"ur Kartographie und Geod\"asie, Frankfurt am Main.
\bibitem{sd81} Smith, M.L.\ \& Dahlen, F.A., 1981. The period and Q of the Chandler wobble,
\emph{Geophys.\ J.\ R.\ astr.\ Soc.}, 64, 223--281.
\bibitem{wahr82} Wahr, J.M., 1982. The effects of the atmosphere and oceans on the Earth's wobble --- I.
Theory, \emph{Geophys.\ J.\ R.\ astr.\ Soc.}, 70, 349--372.
\end{thebibliography}

\end{document}
